\chapter{Logica, verzamelingen en afbeeldingen}\label{ch:b1:logic}

Tot nu toe steunden bewijzen op een informeel, maar eerlijk gevoel voor
wat ``bewijzen'' inhoudt. Dit eerste hoofdstuk van de universitaire
wiskunde legt de spelregels vast: wat een wiskundige \omterm{def:b1:logic:statement}{uitspraak} is, hoe
connectieven en kwantoren \omterm{def:b1:logic:statement}{uitspraken} aan elkaar knopen, welke zetten in
een bewijs geoorloofd zijn --- en bouwt op dat fundament de twee
universele talen van de wiskunde: \omterm{def:b1:logic:sets}{verzamelingen} en \omterm{def:b1:logic:map}{afbeeldingen}.

\section{Uitspraken en connectieven}

\begin{definition}[Uitspraak, connectieven]\label{def:b1:logic:statement}
Een \emph{uitspraak}\index{uitspraak} (of \emph{propositie}) is een zin
die waar (W) of onwaar (O) is --- precies één van beide. Uit twee
uitspraken $P$ en $Q$ vormen we:
\begin{itemize}
  \item de \emph{negatie} $\lnot P$ (``niet $P$''), waar precies wanneer
        $P$ onwaar is;
  \item de \emph{conjunctie} $P \land Q$ (``$P$ en $Q$''), waar precies
        wanneer beide waar zijn;
  \item de \emph{disjunctie} $P \lor Q$ (``$P$ of $Q$''), waar precies
        wanneer minstens één van beide waar is (deze ``of'' is
        inclusief);
  \item de \emph{implicatie} $P \implies Q$, onwaar precies wanneer $P$
        waar is en $Q$ onwaar;
  \item de \emph{equivalentie} $P \iff Q$, waar precies wanneer $P$ en
        $Q$ dezelfde waarheidswaarde hebben.
\end{itemize}
\end{definition}

\begin{remark}
Bij de waarheidstabel van $P \implies Q$ loont het even stil te staan:
zodra $P$ onwaar is, is $P \implies Q$ \emph{waar}, wat $Q$ ook is. ``Als
$2 < 1$, dan $0 = 5$'' is een ware implicatie. Een implicatie beweert
niets over het geval waarin haar hypothese niet opgaat.
\end{remark}

\begin{proposition}[Rekenregels voor uitspraken]\label{prop:b1:logic:rules}
Voor alle \omterm{def:b1:logic:statement}{uitspraken} $P$, $Q$, $R$ geldt:
\begin{enumerate}
  \item $\lnot(\lnot P) \iff P$;
  \item de wetten van De Morgan\index{wetten van De Morgan}:
        $\lnot(P \land Q) \iff (\lnot P) \lor (\lnot Q)$ en
        $\lnot(P \lor Q) \iff (\lnot P) \land (\lnot Q)$;
  \item $(P \implies Q) \iff \bigl((\lnot P) \lor Q\bigr)$, en dus
        $\lnot(P \implies Q) \iff P \land (\lnot Q)$;
  \item contrapositie\index{contrapositie}:
        $(P \implies Q) \iff \bigl((\lnot Q) \implies (\lnot P)\bigr)$;
  \item $(P \iff Q) \iff \bigl((P \implies Q) \land (Q \implies
        P)\bigr)$;
  \item distributiviteit: $P \land (Q \lor R) \iff (P \land Q) \lor
        (P \land R)$ en $P \lor (Q \land R) \iff (P \lor Q) \land
        (P \lor R)$.
\end{enumerate}
\end{proposition}

\begin{proof}
Elke equivalentie controleer je door waarheidstabellen te vergelijken:
twee samengestelde \omterm{def:b1:logic:statement}{uitspraken} in $P$, $Q$, $R$ zijn equivalent precies
wanneer ze in elk van de (vier of acht) gevallen dezelfde waarheidswaarde
aannemen. We schrijven één tabel volledig uit, voor de eerste wet van De
Morgan:
\begin{center}
\begin{tabular}{cc|c|c|cc|c}
$P$ & $Q$ & $P \land Q$ & $\lnot(P \land Q)$ & $\lnot P$ & $\lnot Q$ &
$(\lnot P) \lor (\lnot Q)$ \\
\hline
W & W & W & O & O & O & O \\
W & O & O & W & O & W & W \\
O & W & O & W & W & O & W \\
O & O & O & W & W & W & W \\
\end{tabular}
\end{center}
Kolom $4$ en kolom $7$ vallen samen, en dat bewijst de wet. Voor de
contrapositie gaat het sneller met woorden: $P \implies Q$ is onwaar
precies in het geval ($P$ waar, $Q$ onwaar), en $(\lnot Q) \implies
(\lnot P)$ is onwaar precies in het geval ($\lnot Q$ waar, $\lnot P$
onwaar), oftewel ($Q$ onwaar, $P$ waar) --- hetzelfde enige geval, dus
hebben beide implicaties identieke tabellen. De overige regels gaan
net zo. Merk op dat (3) elke implicatie herleidt tot een disjunctie,
zodat (2) er werktuiglijk de negatieregel $\lnot(P \implies Q) \iff P
\land (\lnot Q)$ uit maakt: wie een implicatie wil weerleggen, moet een
geval aanwijzen waarin de hypothese geldt en de conclusie faalt.
\end{proof}

\section{Kwantoren}

\begin{definition}[Kwantoren]\label{def:b1:logic:quantifiers}
Zij $P(x)$ een eigenschap van een element $x$ van een \omterm{def:b1:logic:sets}{verzameling} $E$.
\begin{itemize}
  \item $\forall x \in E,\ P(x)$ (``voor alle $x$ in $E$ geldt $P(x)$'')
        is waar wanneer elk element van $E$ aan $P$ voldoet;
  \item $\exists x \in E,\ P(x)$ (``er bestaat een $x$ in $E$ zodat
        $P(x)$'') is waar wanneer minstens één element van $E$ aan $P$
        voldoet.
\end{itemize}
We schrijven $\exists!$ voor ``er bestaat precies één''.
\end{definition}

\begin{proposition}[Negatie van kwantoren]\label{prop:b1:logic:negquant}
\[
\lnot\bigl(\forall x \in E,\ P(x)\bigr) \iff
\exists x \in E,\ \lnot P(x),
\qquad
\lnot\bigl(\exists x \in E,\ P(x)\bigr) \iff
\forall x \in E,\ \lnot P(x).
\]
\end{proposition}

\begin{proof}
We bewijzen de eerste equivalentie in beide richtingen; de tweede gaat
symmetrisch. Is $\forall x \in E,\ P(x)$ onwaar, dan voldoet niet elk
element aan $P$: de \omterm{def:b1:logic:sets}{verzameling} $A = \{x \in E : \lnot P(x)\}$ kan niet
leeg zijn, en elk van haar elementen getuigt van $\exists x \in E,\
\lnot P(x)$. Omgekeerd, als een zekere $x_0 \in E$ voldoet aan $\lnot
P(x_0)$, dan is $x_0$ een tegenvoorbeeld en faalt de universele
\omterm{def:b1:logic:statement}{uitspraak}. Voor de tweede regel: ``geen enkele $x$ voldoet aan $P$''
zegt dat de \omterm{def:b1:logic:sets}{verzameling} $\{x : P(x)\}$ leeg is, dat wil zeggen dat elke
$x$ in haar complement $A$ ligt. Passen we beide regels in cascade toe
op een rij geneste kwantoren, dan krijgen we de mechanische procedure
van \cref{ex:b1:logic:limit}: de negatie loopt van links naar rechts,
verwisselt elke $\forall$ met $\exists$ en elke $\exists$ met $\forall$,
en ontkent ten slotte het binnenste predicaat.
\end{proof}

\begin{example}[Alledaagse wiskundige zinnen ontkennen]\label{ex:b1:logic:negpractice}
Zij $f \colon \R \to \R$. De zin ``$f$ is stijgend'' luidt
\[
\forall x \in \R,\ \forall y \in \R,\quad
x \leq y \implies f(x) \leq f(y) ,
\]
en haar negatie is, volgens \cref{prop:b1:logic:negquant} samen met de
regel $\lnot(P \implies Q) \iff P \land \lnot Q$:
\[
\exists x \in \R,\ \exists y \in \R,\quad
x \leq y \ \text{ en }\ f(x) > f(y) :
\]
één getuigend paar volstaat. Evenzo is ``$f$ is begrensd'' de \omterm{def:b1:logic:statement}{uitspraak}
$\exists M \in \R,\ \forall x \in \R,\ \abs{f(x)} \leq M$, met negatie
\[
\forall M \in \R,\ \exists x \in \R,\quad \abs{f(x)} > M :
\]
welke grens ook wordt voorgesteld, er is een punt dat haar overtreft.
Het inzicht: een correcte negatie bevat nooit een ``niet'' vóór een blok
kwantoren --- ze is een nieuwe, positieve \omterm{def:b1:logic:statement}{uitspraak} waarin de rollen
omgedraaid zijn: je levert nu de getuigen die je eerst kreeg aangereikt.
\end{example}

\begin{example}[Volgorde van de kwantoren]\label{ex:b1:logic:order}
De volgorde van verschillende kwantoren doet ertoe:
\[
\forall x \in \R,\ \exists y \in \R,\ y > x
\quad\text{is waar (neem } y = x+1\text{),}
\]
\[
\exists y \in \R,\ \forall x \in \R,\ y > x
\quad\text{is onwaar (geen reëel getal overtreft alle reële getallen).}
\]
In de eerste \omterm{def:b1:logic:statement}{uitspraak} mag $y$ van $x$ afhangen; in de tweede moet één
enkele $y$ het voor alle $x$ doen. Twee gelijke kwantoren mogen daarentegen
altijd van plaats wisselen.
\end{example}

\begin{example}[Een definitie met drie kwantoren lezen]\label{ex:b1:logic:limit}
De zin ``de rij $(u_n)$ convergeert naar $\ell$'' wordt in
\cref{ch:b1:seq} geschreven als
\[
\forall \varepsilon > 0,\ \exists N \in \N,\ \forall n \geq N,\quad
\abs{u_n - \ell} \leq \varepsilon .
\]
Haar negatie luidt, na driemaal \cref{prop:b1:logic:negquant},
\[
\exists \varepsilon > 0,\ \forall N \in \N,\ \exists n \geq N,\quad
\abs{u_n - \ell} > \varepsilon .
\]
Zulke zinnen werktuiglijk kunnen ontkennen, zonder na te denken over wat
ze betekenen, is een vaardigheid op zich: ze scheidt het logische werk
van het wiskundige.
\end{example}

\section{Bewijstechnieken}

\begin{method}[De standaardpatronen van een bewijs]\label{met:b1:logic:proofs}
Om te bewijzen\dots
\begin{enumerate}
  \item \emph{een implicatie $P \implies Q$ rechtstreeks}: neem $P$ aan,
        leid $Q$ af;
  \item \emph{via contrapositie}: neem $\lnot Q$ aan, leid $\lnot P$ af
        --- geldig volgens \cref{prop:b1:logic:rules}~(4);
  \item \emph{uit het ongerijmde}: neem aan dat de \omterm{def:b1:logic:statement}{uitspraak} onwaar is en
        leid een tegenspraak af;
  \item \emph{een equivalentie}: bewijs beide implicaties apart (of rijg
        bekende equivalenties aaneen);
  \item \emph{een ``voor alle''-uitspraak}: kies een \emph{willekeurige}
        $x$ in $E$ (``zij $x \in E$'') en bewijs $P(x)$;
  \item \emph{een ``er bestaat''-uitspraak}: wijs een getuige aan, of
        bewijs het bestaan langs een omweg;
  \item \emph{met inductie}: zie \cref{thm:b1:logic:induction}.
\end{enumerate}
Bewijs je iets over een goedgekozen maar willekeurig element, geef dat
element dan nooit extra eigenschappen: ``zij $x \in \R$'' gevolgd door
``omdat $x > 0$\dots'' bewijst niets over negatieve $x$.
\end{method}

\begin{remark}[Veelgemaakte fouten in bewijzen]\label{rem:b1:logic:pitfalls}
Vier klassieke valkuilen, die het verdienen één keer bij naam genoemd te
worden.
\begin{enumerate}
  \item \emph{Omkering in plaats van contrapositie.} $Q \implies P$ is
        \emph{niet} equivalent met $P \implies Q$; alleen $\lnot Q
        \implies \lnot P$ is dat. ``Als het regent, is de straat nat''
        geeft je niet het recht uit een natte straat regen te besluiten.
  \item \emph{Een equivalentie met één implicatie bewijzen.} Een
        ``dan en slechts dan''-bewering is twee stellingen; zeg welke
        richting je bewijst, en bewijs ze allebei. Een keten van $\iff$
        is alleen geldig als \emph{elke} schakel werkelijk omkeerbaar is
        --- een vergelijking kwadrateren bijvoorbeeld is dat niet.
  \item \emph{Achterstevoren bewijzen.} Vertrekken van de gewenste
        conclusie en daaruit iets waars afleiden bewijst niets (uit
        $-1 = 1$ volgt na kwadrateren het ware $1 = 1$). Een berekening
        mag achterstevoren \emph{gevonden} worden, maar ze moet
        voorwaarts \emph{opgeschreven} worden, of met expliciete
        equivalenties.
  \item \emph{Vaste getuige tegenover willekeurig element.} Voor
        $\exists x,\ P(x)$ volstaat één slim gekozen $x$; voor
        $\forall x,\ P(x)$ moet de gekozen $x$ willekeurig blijven. Die
        twee door elkaar halen --- een universele bewering op een
        voorbeeld nagaan --- is de meest voorkomende fout in
        beginnerswerk.
\end{enumerate}
\end{remark}

\begin{example}[Contrapositie en ongerijmde aan het werk]\label{ex:b1:logic:sqrt2}
\emph{Voor $n \in \N$: is $n^2$ even, dan is $n$ even.} Via
contrapositie: is $n$ oneven, $n = 2k+1$, dan is $n^2 = 4k^2 + 4k + 1$
oneven.

\emph{$\sqrt 2$ is irrationaal.} Uit het ongerijmde: stel $\sqrt 2 = p/q$
met $p, q \in \N^*$ en de breuk onvereenvoudigbaar. Dan is $p^2 = 2q^2$
even, dus is $p$ even (vorig punt), zeg $p = 2r$; dan is $q^2 = 2r^2$
even, dus is $q$ even --- in tegenspraak met de onvereenvoudigbaarheid.
\end{example}

\begin{theorem}[Inductie]\label{thm:b1:logic:induction}
Zij $P(n)$ een eigenschap van het gehele getal $n$. Als
\begin{enumerate}
  \item $P(0)$ waar is, en
  \item voor alle $n \in \N$ geldt $P(n) \implies P(n+1)$,
\end{enumerate}
dan is $P(n)$ waar voor alle $n \in \N$.

\emph{Sterke inductie:} dezelfde conclusie geldt wanneer (2) vervangen
wordt door: voor alle $n$ geldt $\bigl(P(0) \land \dots \land
P(n)\bigr) \implies P(n+1)$.
\end{theorem}

\begin{proof}
Dit is een eigenschap van $\N$ zelf, equivalent met: \emph{elke
niet-lege deelverzameling van $\N$ heeft een kleinste element} (wat we
als bekend aannemen). Stel immers dat (1) en (2) gelden en zet $A = \{n
\in \N : P(n) \text{ onwaar}\}$. Is $A \neq \emptyset$, dan heeft $A$
een kleinste element $m$; wegens (1) is $m \neq 0$; dan is $m - 1 \notin
A$, dus geldt $P(m-1)$, en (2) levert $P(m)$ --- tegenspraak. Dus $A =
\emptyset$. Voor sterke inductie werkt hetzelfde argument: $P(0), \dots,
P(m-1)$ gelden alle, omdat $m$ het kleinste element van $A$ is.
\end{proof}

\begin{example}[Uniek bestaan bewijzen]\label{ex:b1:logic:uniqueexistence}
Een \omterm{def:b1:logic:statement}{uitspraak} $\exists!\,x,\ P(x)$ bestaat uit \emph{twee} \omterm{def:b1:logic:statement}{uitspraken},
die apart bewezen worden: het bestaan (wijs een $x_0$ met $P(x_0)$ aan of
construeer er een) en de uniciteit (neem $P(x)$ en $P(x')$ aan en leid
$x = x'$ af). Een voorbeeld: \emph{er is precies één reëel getal $x$ met
$x^3 + x = 2$.} Bestaan: $x_0 = 1$ voldoet, want $1 + 1 = 2$. Uniciteit:
uit $x^3 + x = x'^3 + x'$ volgt
\[
0 = (x^3 - x'^3) + (x - x')
= (x - x')\,\bigl(x^2 + xx' + x'^2 + 1\bigr),
\]
en de tweede factor is positief (hij is gelijk aan $\bigl(x +
\tfrac{x'}2\bigr)^2 + \tfrac34 x'^2 + 1 \geq 1$), zodat $x = x'$. Let op
de werkverdeling: het bestaan berustte op een gelukkige gok, de uniciteit
op algebra die geldt voor \emph{willekeurige} oplossingen --- geen van
beide argumenten doet het werk van het andere, en de tweede helft
vergeten blijft een blijvende verleiding zodra één oplossing gevonden is.
\end{example}

\begin{example}\label{ex:b1:logic:induction}
Voor alle $n \in \N^*$ geldt $\;\sum_{k=1}^n k = \frac{n(n+1)}{2}$.
Basisgeval $n = 1$: beide leden zijn $1$. Stap: neem de formule aan voor
$n$, dan is
\[
\sum_{k=1}^{n+1} k = \frac{n(n+1)}{2} + (n+1)
= (n+1)\Bigl(\frac n2 + 1\Bigr) = \frac{(n+1)(n+2)}{2}. \qedhere
\]
\end{example}

\begin{example}[Sterke inductie aan het werk]\label{ex:b1:logic:strongind}
\emph{Elk geheel getal $n \geq 2$ is een product van priemgetallen} (een
priemgetal is een geheel getal $\geq 2$ waarvan $1$ en het getal zelf de
enige delers $\geq 1$ zijn; priemgetallen worden om hun eigen wille
bestudeerd in \cref{ch:b1:arith}). Gewone inductie helpt hier niet:
weten dat $95 = 5 \times 19$ ontbindt, zegt niets over $96$. Sterke
inductie past precies. Basisgeval: $2$ is priem, dus een product van
priemgetallen met één factor. Stap: zij $n \geq 2$ en neem aan dat elk
geheel getal $m$ met $2 \leq m \leq n$ een product van priemgetallen is.
Is $n + 1$ priem, dan zijn we klaar. Zo niet, dan is $n + 1 = ab$ met
$2 \leq a, b \leq n$; volgens de sterke hypothese zijn $a$ en $b$
producten van priemgetallen, en dus ook $n + 1$. Het inzicht: sterke
inductie is het juiste gereedschap zodra de ``reden'' voor $P(n+1)$ op een
onvoorspelbare eerdere rang ligt en niet op rang $n$.
\end{example}

\section{Verzamelingen}

\begin{definition}[Bewerkingen met verzamelingen]\label{def:b1:logic:sets}
We nemen het begrip \emph{verzameling}\index{verzameling} en de
lidmaatschapsrelatie $x \in E$ als primitief aan. Voor verzamelingen $A,
B$ binnen een omvattende verzameling $E$:
\begin{itemize}
  \item \emph{inclusie}: $A \subseteq B$ wanneer $\forall x,\ x \in A
        \implies x \in B$; gelijkheid $A = B$ wanneer $A \subseteq B$ en
        $B \subseteq A$;
  \item de \emph{vereniging} $A \cup B$, de \emph{doorsnede} $A \cap B$,
        het \emph{verschil} $A \setminus B = \{x \in A : x \notin B\}$,
        het \emph{complement} $\overline{A} = E \setminus A$;
  \item de \emph{lege verzameling} $\emptyset$, bevat in elke
        verzameling;
  \item de \emph{machtsverzameling}\index{machtsverzameling}
        $\mathcal{P}(E)$: de verzameling van alle deelverzamelingen van
        $E$;
  \item het \emph{product} $E \times F$: de verzameling van geordende
        paren $(x, y)$ met $x \in E$, $y \in F$.
\end{itemize}
\end{definition}

\begin{example}[Wennen aan de machtsverzameling]\label{ex:b1:logic:powerset}
Voor $E = \{a, b\}$ is
\[
\mathcal P(E) = \bigl\{\, \emptyset,\ \{a\},\ \{b\},\ \{a, b\}
\,\bigr\},
\]
vier elementen --- en let op de typediscipline: $a \in E$, maar $\{a\}
\in \mathcal P(E)$; de \omterm{def:b1:logic:statement}{uitspraken} $a \in \mathcal P(E)$ en $\{a\}
\subseteq \mathcal P(E)$ zijn zoals ze er staan allebei onwaar (de tweede
zou vergen dat $a$ een \emph{deelverzameling} van $E$ is). Iteratie
vanuit het niets: $\mathcal P(\emptyset) = \{\emptyset\}$ heeft één
element, $\mathcal P(\mathcal P(\emptyset)) = \{\emptyset,
\{\emptyset\}\}$ er twee, de volgende vier --- \omterm{def:b1:logic:sets}{verzamelingen} van
\omterm{def:b1:logic:sets}{verzamelingen} zijn gewone \omterm{def:b1:logic:sets}{verzamelingen}, en \cref{ch:b1:counting} zal het
verdubbelingspatroon bevestigen: $\abs{\mathcal P(E)} = 2^{\abs E}$. De
niveaus ($x$, $\{x\}$, $\{\{x\}\}$) uit elkaar houden is het halve werk
bij oefeningen als \cref{exo:b1:logic:11,exo:b1:logic:12}.
\end{example}

\begin{proposition}[Verzamelingenalgebra]\label{prop:b1:logic:setalgebra}
Voor deelverzamelingen $A, B, C$ van $E$ geldt:
\begin{enumerate}
  \item $A \cap (B \cup C) = (A \cap B) \cup (A \cap C)$ en
        $A \cup (B \cap C) = (A \cup B) \cap (A \cup C)$;
  \item De Morgan: $\overline{A \cup B} = \overline{A} \cap
        \overline{B}$ en $\overline{A \cap B} = \overline{A} \cup
        \overline{B}$;
  \item $A \subseteq B \iff \overline{B} \subseteq \overline{A}$.
\end{enumerate}
\end{proposition}

\begin{proof}
Elke identiteit vertaalt een regel uit \cref{prop:b1:logic:rules} via
het woordenboek (wel of niet $\in A$) $\leftrightarrow$ (\omterm{def:b1:logic:statement}{uitspraak} waar
of onwaar): bijvoorbeeld $x \in \overline{A \cup B} \iff \lnot(x \in A
\lor x \in B) \iff (x \notin A) \land (x \notin B) \iff x \in
\overline{A} \cap \overline{B}$. Punt (3) is contrapositie. Als tweede
staaltje de eerste distributiviteitswet, volledig uitgeschreven:
\[
x \in A \cap (B \cup C)
\iff (x \in A) \land \bigl(x \in B \lor x \in C\bigr)
\iff \bigl(x \in A \land x \in B\bigr) \lor
\bigl(x \in A \land x \in C\bigr),
\]
wegens de distributiviteit uit \cref{prop:b1:logic:rules}~(6), en de
laatste \omterm{def:b1:logic:statement}{uitspraak} zegt $x \in (A \cap B) \cup (A \cap C)$. Elke
verzamelingsidentiteit van dit type volgt uit diezelfde ene mechanische
vertaling --- en daarom hoeft er geen enkele uit het hoofd geleerd te
worden.
\end{proof}

\begin{method}[Gelijkheid van verzamelingen bewijzen]\label{met:b1:logic:doubleinc}
Om $A = B$ te bewijzen, bewijs je de twee inclusies: zij $x \in A$, toon
$x \in B$; zij vervolgens $x \in B$, toon $x \in A$. Als alternatief rijg
je equivalenties aaneen, $x \in A \iff \dots \iff x \in B$, mits elke
stap werkelijk een equivalentie is.
\end{method}

\begin{omfigure}
\begin{tikzpicture}[scale=0.9]
  \draw[thick] (-2.6,-1.9) rectangle (2.6,1.9);
  \node[below right] at (-2.6,1.9) {$E$};
  \begin{scope}
    \clip (-2.6,-1.9) rectangle (2.6,1.9);
    \fill[omDef!20] (-2.6,-1.9) rectangle (2.6,1.9);
    \fill[white] (-0.7,0) circle (1.1);
    \fill[white] (0.7,0) circle (1.1);
  \end{scope}
  \draw[omDef, thick] (-0.7,0) circle (1.1);
  \draw[omDef, thick] (0.7,0) circle (1.1);
  \node[omDef] at (-1.35,0.5) {$A$};
  \node[omDef] at (1.35,0.5) {$B$};
  \begin{scope}[xshift=6.6cm]
    \draw[thick] (-2.6,-1.9) rectangle (2.6,1.9);
    \node[below right] at (-2.6,1.9) {$E$};
    \begin{scope}
      \clip (-2.6,-1.9) rectangle (2.6,1.9);
      \fill[omThm!20] (-2.6,-1.9) rectangle (2.6,1.9);
      \begin{scope}
        \clip (-0.7,0) circle (1.1);
        \fill[white] (0.7,0) circle (1.1);
      \end{scope}
    \end{scope}
    \draw[omThm, thick] (-0.7,0) circle (1.1);
    \draw[omThm, thick] (0.7,0) circle (1.1);
    \node[omThm] at (-1.35,0.5) {$A$};
    \node[omThm] at (1.35,0.5) {$B$};
  \end{scope}
\end{tikzpicture}

{\small De wetten van De Morgan in beeld: het gearceerde gebied links is
$\overline{A \cup B} = \overline A \cap \overline B$ (alles buiten beide
schijven); rechts is het $\overline{A \cap B} = \overline A \cup
\overline B$ (alles behalve de lensvormige overlap). Een tekening is geen
bewijs, maar ze maakt het elementsgewijze bewijs van
\cref{prop:b1:logic:setalgebra} onvergetelijk.}
\end{omfigure}

\section{Afbeeldingen}

\begin{definition}[Afbeelding, beeld, origineel]\label{def:b1:logic:map}
Een \emph{afbeelding}\index{afbeelding} (of \emph{functie}) $f \colon E
\to F$ kent aan elk element $x$ van de \omterm{def:b1:logic:sets}{verzameling} $E$ (het
\emph{domein}) precies één element $f(x)$ van de \omterm{def:b1:logic:sets}{verzameling} $F$ (het
\emph{codomein}) toe. Voor $A \subseteq E$ en $B \subseteq F$ zijn
\[
f(A) = \{f(x) : x \in A\} \subseteq F,
\qquad
f^{-1}(B) = \{x \in E : f(x) \in B\} \subseteq E
\]
het \emph{beeld} van $A$ en het \emph{origineel}\index{origineel} van
$B$. De \emph{samenstelling} van $f \colon E \to F$ en $g \colon F \to G$
is $g \circ f \colon E \to G$, $x \mapsto g(f(x))$.
\end{definition}

\begin{remark}
De notatie $f^{-1}(B)$ veronderstelt \emph{geen} inverse \omterm{def:b1:logic:map}{afbeelding}:
$f^{-1}(B)$ is voor elke $f$ gedefinieerd. \omterm{def:b1:logic:map}{Originelen} gedragen zich beter
dan beelden: $f^{-1}$ bewaart verenigingen, doorsneden en complementen,
terwijl $f(A \cap A') \subseteq f(A) \cap f(A')$ strikt kan zijn
(\cref{exo:b1:logic:8}).
\end{remark}

\begin{example}[Beelden en originelen berekenen]\label{ex:b1:logic:images}
Zij $f \colon \R \to \R$, $x \mapsto x^2$. Dan is
\[
f\bigl(\intcc{-1}{2}\bigr) = \intcc04, \qquad
f^{-1}\bigl(\intcc14\bigr) = \intcc{-2}{-1} \cup \intcc12, \qquad
f^{-1}(\{-1\}) = \emptyset .
\]
Voor het eerste: elke $x \in \intcc{-1}2$ heeft $x^2 \in \intcc04$, en
elke $y \in \intcc04$ wordt bereikt als $y = (\sqrt y)^2$ met $\sqrt y
\in \intcc02 \subseteq \intcc{-1}2$ --- merk op dat het beeld \emph{niet}
$\intcc14 = \{(-1)^2, 2^2\}$ is: beelden van intervallen bereken je niet
uit de randpunten alleen. Voor het tweede: $1 \leq x^2 \leq 4 \iff 1 \leq
\abs x \leq 2$, wat in twee stukken uiteenvalt. Het derde toont dat een
\omterm{def:b1:logic:map}{origineel} leeg mag zijn --- $f^{-1}(B)$ is altijd zinvol, hoe klein de
doorsnede van $B$ met het beeld ook is. Bekijk ten slotte op ditzelfde
voorbeeld het strengheidsverschijnsel uit de opmerking hierboven: met $A
= \intcc{-1}0$ en $A' = \intcc01$ is $f(A \cap A') = f(\{0\}) = \{0\}$,
terwijl $f(A) \cap f(A') = \intcc01$.
\end{example}

\begin{definition}[Injectief, surjectief, bijectief]\label{def:b1:logic:inj}
Een \omterm{def:b1:logic:map}{afbeelding} $f \colon E \to F$ heet:
\begin{itemize}
  \item \emph{injectief}\index{injectief} wanneer verschillende elementen
        verschillende beelden hebben: $\forall x, x' \in E,\ f(x) = f(x')
        \implies x = x'$;
  \item \emph{surjectief}\index{surjectief} wanneer elk element van $F$
        bereikt wordt: $\forall y \in F,\ \exists x \in E,\ f(x) = y$;
  \item \emph{bijectief}\index{bijectief} wanneer ze beide is, dat wil
        zeggen wanneer elke $y \in F$ precies één \omterm{def:b1:logic:map}{origineel} heeft.
\end{itemize}
\end{definition}

\begin{theorem}[Inverse afbeelding]\label{thm:b1:logic:inverse}
Een \omterm{def:b1:logic:map}{afbeelding} $f \colon E \to F$ is \omterm{def:b1:logic:inj}{bijectief} dan en slechts dan als er
een \omterm{def:b1:logic:map}{afbeelding} $g \colon F \to E$ bestaat met $g \circ f =
\mathrm{id}_E$ en $f \circ g = \mathrm{id}_F$. In dat geval is $g$ uniek;
ze wordt genoteerd $f^{-1}$ en de \emph{inverse} van $f$ genoemd, en
$f^{-1}$ is zelf \omterm{def:b1:logic:inj}{bijectief} met $(f^{-1})^{-1} = f$.
\end{theorem}

\begin{proof}
($\Rightarrow$) Is $f$ \omterm{def:b1:logic:inj}{bijectief}, dan heeft elke $y \in F$ precies één
\omterm{def:b1:logic:map}{origineel}; definieer $g(y)$ als dat \omterm{def:b1:logic:map}{origineel}. Dan is $f(g(y)) = y$ per
constructie, en $g(f(x)) = x$ omdat $x$ hét \omterm{def:b1:logic:map}{origineel} van $f(x)$ is.

($\Leftarrow$) Stel dat zo'n $g$ bestaat. Uit $f(x) = f(x')$ volgt na
toepassing van $g$ dat $x = x'$: $f$ is \omterm{def:b1:logic:inj}{injectief}. Voor $y \in F$
voldoet $x = g(y)$ aan $f(x) = y$: $f$ is \omterm{def:b1:logic:inj}{surjectief}.

Uniciteit: voldoen zowel $g$ als $h$, dan is $g = g \circ \mathrm{id}_F
= g \circ (f \circ h) = (g \circ f) \circ h = h$. Ten slotte is het paar
identiteiten symmetrisch in $f$ en $g$, zodat $g = f^{-1}$ \omterm{def:b1:logic:inj}{bijectief} is
met inverse $f$.
\end{proof}

\begin{example}[Een inverse in de praktijk berekenen]\label{ex:b1:logic:inversecomp}
Zij $f \colon \R \to \intoo0{+\infty}$, $f(x) = \eu^{2x+1}$. Om te
inverteren los je $y = f(x)$ op naar $x$, voor gegeven $y > 0$:
\[
y = \eu^{2x+1} \iff \ln y = 2x + 1 \iff x = \frac{\ln y - 1}2 ,
\]
en elke stap is omkeerbaar op de aangekondigde domeinen. De berekening
levert alles tegelijk: bij elke $y$ in het codomein hoort precies één
oplossing $x$, dus $f$ is \omterm{def:b1:logic:inj}{bijectief}, en
\[
f^{-1} \colon \intoo0{+\infty} \to \R,
\qquad
f^{-1}(y) = \frac{\ln y - 1}2 .
\]
Een snelle controle van beide samenstellingen ($f^{-1}(f(x)) =
\frac{(2x+1) - 1}2 = x$ en $f(f^{-1}(y)) = \eu^{\ln y} = y$) bevestigt
het criterium van \cref{thm:b1:logic:inverse}. Het inzicht: ``los op naar
$x$ en houd de equivalenties in het oog'' is tegelijk het bestaansbewijs,
het uniciteitsbewijs en de formule --- maar het werkt alleen als het
codomein correct is aangekondigd ($f$ is \emph{niet} \omterm{def:b1:logic:inj}{surjectief} op $\R$).
\end{example}

\begin{proposition}[Samenstelling en de drie eigenschappen]\label{prop:b1:logic:comp}
Zij $f \colon E \to F$ en $g \colon F \to G$.
\begin{enumerate}
  \item Zijn $f$ en $g$ \omterm{def:b1:logic:inj}{injectief} (respectievelijk \omterm{def:b1:logic:inj}{surjectief},
        \omterm{def:b1:logic:inj}{bijectief}), dan is $g \circ f$ dat ook; en in het \omterm{def:b1:logic:inj}{bijectieve}
        geval is $(g \circ f)^{-1} = f^{-1} \circ g^{-1}$.
  \item Is $g \circ f$ \omterm{def:b1:logic:inj}{injectief}, dan is $f$ \omterm{def:b1:logic:inj}{injectief}. Is $g \circ f$
        \omterm{def:b1:logic:inj}{surjectief}, dan is $g$ \omterm{def:b1:logic:inj}{surjectief}.
\end{enumerate}
\end{proposition}

\begin{proof}
(1) Uit $g(f(x)) = g(f(x'))$ volgt met de injectiviteit van $g$ dat
$f(x) = f(x')$, en met die van $f$ dat $x = x'$. Is $z \in G$, dan levert
de surjectiviteit van $g$ een $y$ met $g(y) = z$, en die van $f$ een $x$
met $f(x) = y$, zodat $g(f(x)) = z$. In het \omterm{def:b1:logic:inj}{bijectieve} geval ga je
rechtstreeks na dat $f^{-1} \circ g^{-1}$ een tweezijdige inverse van $g
\circ f$ is, waarna de uniciteit in \cref{thm:b1:logic:inverse} de zaak
afmaakt.

(2) Uit $f(x) = f(x')$ volgt $g(f(x)) = g(f(x'))$, en de injectiviteit
van $g \circ f$ geeft $x = x'$. Is $z \in G$, dan levert de surjectiviteit
van $g \circ f$ een $x$ met $g(f(x)) = z$: dan voldoet $y = f(x)$ aan
$g(y) = z$.
\end{proof}

\begin{example}[Punt (2) is scherp]\label{ex:b1:logic:compsharp}
In \cref{prop:b1:logic:comp}~(2) kunnen de conclusies niet versterkt
worden: dat $g \circ f$ \omterm{def:b1:logic:inj}{bijectief} is, dwingt $f$ \emph{niet} \omterm{def:b1:logic:inj}{surjectief}
of $g$ \omterm{def:b1:logic:inj}{injectief} te zijn. Neem $E = G = \{1\}$, $F = \{1, 2\}$, met
$f(1) = 1$ en $g(1) = g(2) = 1$: dan is $g \circ f = \mathrm{id}_E$
\omterm{def:b1:logic:inj}{bijectief}, terwijl $f$ het element $2$ mist en $g$ beide elementen op
elkaar plakt. De moraal is een nauwkeurige boekhoudregel: informatie over
de samenstelling stroomt voor injectiviteit naar de \emph{binnenste}
\omterm{def:b1:logic:map}{afbeelding} en voor surjectiviteit naar de \emph{buitenste}, nooit
andersom. (\cref{exo:b1:logic:9} bouwt hetzelfde verschijnsel op oneindige
\omterm{def:b1:logic:sets}{verzamelingen}, waar het de motor is achter eenzijdige inversen.)
\end{example}

\begin{example}\label{ex:b1:logic:injexamples}
$f \colon \R \to \R$, $x \mapsto x^2$ is \omterm{def:b1:logic:inj}{injectief} noch \omterm{def:b1:logic:inj}{surjectief}
($f(-1) = f(1)$, en $-1$ heeft geen \omterm{def:b1:logic:map}{origineel}). Beperken we domein en
codomein, dan is $f \colon \R_+ \to \R_+$, $x \mapsto x^2$ \omterm{def:b1:logic:inj}{bijectief}, met
inverse $y \mapsto \sqrt y$. Of een \omterm{def:b1:logic:map}{afbeelding} \omterm{def:b1:logic:inj}{injectief} of \omterm{def:b1:logic:inj}{surjectief}
is, hangt dus af van het aangekondigde domein en codomein, niet alleen
van het voorschrift.
\end{example}

\section{Relaties}

\begin{definition}[Equivalentierelatie]\label{def:b1:logic:equiv}
Een \emph{binaire relatie} $\mathcal{R}$ op een \omterm{def:b1:logic:sets}{verzameling} $E$ heet een
\emph{equivalentierelatie}\index{equivalentierelatie} wanneer ze
\emph{reflexief} is ($x \mathbin{\mathcal{R}} x$ voor alle $x$),
\emph{symmetrisch} ($x \mathbin{\mathcal{R}} y \implies y
\mathbin{\mathcal{R}} x$) en \emph{transitief} (uit $x
\mathbin{\mathcal{R}} y$ en $y \mathbin{\mathcal{R}} z$ volgt $x
\mathbin{\mathcal{R}} z$). De \emph{equivalentieklasse} van $x$ is
$\mathrm{cl}(x) = \{y \in E : x \mathbin{\mathcal{R}} y\}$.
\end{definition}

\begin{example}[De drie axioma's nagaan]\label{ex:b1:logic:fracpart}
Verklaar op $\R$ dat $x \mathbin{\mathcal{R}} y$ wanneer $x - y \in \Z$.
\emph{Reflexief:} $x - x = 0 \in \Z$. \emph{Symmetrisch:} is $x - y \in
\Z$, dan is $y - x = -(x - y) \in \Z$. \emph{Transitief:} zijn $x - y \in
\Z$ en $y - z \in \Z$, dan is $x - z = (x - y) + (y - z) \in \Z$ (een som
van gehele getallen). Dus $\mathcal R$ is een \omterm{def:b1:logic:equiv}{equivalentierelatie}, en
$\mathrm{cl}(x) = x + \Z = \{x + k : k \in \Z\}$: elke klasse bevat
precies één vertegenwoordiger in $\intco01$, haar \emph{fractioneel
deel}. De relatie ``$\abs{x - y} \leq 1$'' op $\R$ daarentegen is wel
reflexief en symmetrisch, maar \emph{niet} transitief ($0
\mathbin{\mathcal R} 1$ en $1 \mathbin{\mathcal R} 2$, terwijl $\abs{0 -
2} > 1$): nabijheid plant zich niet voort, en een \omterm{thm:b1:logic:partition}{partitie} in klassen
bestaat niet --- een nuttig tegenvoorbeeld om bij de hand te houden
wanneer het nagaan van de axioma's routine begint te lijken.
\end{example}

\begin{theorem}[Klassen vormen een partitie]\label{thm:b1:logic:partition}
Zij $\mathcal{R}$ een \omterm{def:b1:logic:equiv}{equivalentierelatie} op $E$. Dan zijn de
\omterm{def:b1:logic:equiv}{equivalentieklassen} niet-leeg, twee aan twee disjunct of gelijk, en is
hun vereniging $E$: ze vormen een \emph{partitie}\index{partitie} van
$E$. Omgekeerd komt elke \omterm{thm:b1:logic:partition}{partitie} van $E$ op deze manier van precies één
\omterm{def:b1:logic:equiv}{equivalentierelatie} (``in hetzelfde stuk liggen'').
\end{theorem}

\begin{proof}
Wegens reflexiviteit is $x \in \mathrm{cl}(x)$, dus zijn de klassen
niet-leeg en is hun vereniging $E$. Stel $\mathrm{cl}(x) \cap
\mathrm{cl}(y) \neq \emptyset$, zeg dat $z$ in beide ligt. Dan is $x
\mathbin{\mathcal{R}} z$ en $y \mathbin{\mathcal{R}} z$, dus met symmetrie
en transitiviteit $x \mathbin{\mathcal{R}} y$. Voor elke $t \in
\mathrm{cl}(y)$ geeft transitiviteit nu $t \in \mathrm{cl}(x)$, en
symmetrisch andersom: beide klassen zijn gelijk. Voor de omkering, zij
$(E_i)_{i \in I}$ een \omterm{thm:b1:logic:partition}{partitie} van $E$ en definieer $x \mathbin{\mathcal
S} y$ als ``een zeker stuk bevat zowel $x$ als $y$''. \emph{Reflexief:}
$x$ ligt in een stuk, dat $x$ dan tweemaal bevat. \emph{Symmetrisch:} de
voorwaarde is symmetrisch in $x$ en $y$. \emph{Transitief:} zijn $x, y
\in E_i$ en $y, z \in E_j$, dan is $y \in E_i \cap E_j$, dus $E_i = E_j$
(verschillende stukken zijn disjunct) en delen $x$ en $z$ een stuk. De
$\mathcal S$-klasse van $x$ is precies het stuk dat $x$ bevat, zodat de
klassen de gegeven stukken zijn. Ten slotte ligt de relatie vast door
haar klassen: twee \omterm{def:b1:logic:equiv}{equivalentierelaties} met dezelfde klassen verbinden
dezelfde paren, want elk van beide verbindt $x$ en $y$ precies wanneer
$y$ tot de klasse van $x$ behoort --- waarmee de uniciteit bewezen is.
\end{proof}

\begin{example}\label{ex:b1:logic:congruence}
Op $\Z$ is de congruentie modulo $n$ ($x \equiv y \pmod n$ wanneer $n$
het getal $x - y$ deelt) een \omterm{def:b1:logic:equiv}{equivalentierelatie}; haar klassen zijn de
$n$ \omterm{def:b1:logic:sets}{verzamelingen} van gehele getallen met een gegeven rest bij deling
door $n$. Dit voorbeeld wordt de ring $\Z/n\Z$ in
\cref{ch:b1:structures}.
\end{example}

\begin{definition}[Ordeningsrelatie]\label{def:b1:logic:order}
Een relatie $\preceq$ op $E$ heet een \emph{orde}\index{ordeningsrelatie}
wanneer ze reflexief is, \emph{antisymmetrisch} (uit $x \preceq y$ en $y
\preceq x$ volgt $x = y$) en transitief. De orde is \emph{totaal} wanneer
elke twee elementen vergelijkbaar zijn, en anders \emph{partieel}. Een
element $M \in A \subseteq E$ is een \emph{grootste element} van $A$
wanneer $a \preceq M$ voor alle $a \in A$; grootste (en kleinste)
elementen zijn uniek zodra ze bestaan.
\end{definition}

\begin{example}\label{ex:b1:logic:orders}
$(\R, \leq)$ is totaal geordend. $(\mathcal{P}(E), \subseteq)$ is
partieel geordend zodra $E$ twee elementen heeft: $\{a\}$ en $\{b\}$ zijn
onvergelijkbaar. De deelverzameling $A = \{\{a\}, \{b\}\}$ van
$\mathcal{P}(\{a,b\})$ heeft geen grootste element, maar wel een
bovengrens $\{a, b\}$: het onderscheid tussen grootste elementen en
bovengrenzen keert voor $\R$ terug in \cref{ch:b1:reals}.
\end{example}

\begin{example}[Twee ordes op het rooster $\N^2$]\label{ex:b1:logic:lexorder}
Vergelijk paren natuurlijke getallen componentsgewijs: $(a, b) \preceq
(a', b')$ wanneer $a \leq a'$ \emph{en} $b \leq b'$ (de
\emph{productorde}). Dit is een \omterm{def:b1:logic:order}{orde} --- elk axioma wordt coördinaat voor
coördinaat overgeërfd --- maar een partiële: $(1, 3)$ en $(2, 0)$ zijn
onvergelijkbaar. Vergelijk nu als in een woordenboek: $(a, b)
\preceq_{\mathrm{lex}} (a', b')$ wanneer $a < a'$, of $a = a'$ en $b \leq
b'$ (de \emph{lexicografische \omterm{def:b1:logic:order}{orde}}). Transitiviteit vraagt een
gevalsonderscheid, maar geldt, en nu zijn elke twee paren vergelijkbaar:
de \omterm{def:b1:logic:order}{orde} is totaal. Beide \omterm{def:b1:logic:order}{ordes} rangschikken dezelfde \omterm{def:b1:logic:sets}{verzameling}
verschillend --- $(0, 100) \preceq_{\mathrm{lex}} (1, 0)$ terwijl de
productorde er niets over zegt --- een herinnering dat een \omterm{def:b1:logic:order}{orde} een
structuur is die je \emph{kiest} en geen eigenschap van de \omterm{def:b1:logic:sets}{verzameling}.
Lexicografisch vergelijken is bovendien de standaardtruc om verscheidene
sorteercriteria tot één criterium samen te smeden.
\end{example}

\begin{remark}[Tussenspel: grootte als bijectie]\label{rem:b1:logic:interlude}
Eén stil thema van dit hoofdstuk verdient de schijnwerper: bijecties zijn
het wiskundige begrip van ``even groot''. Voor eindige \omterm{def:b1:logic:sets}{verzamelingen}
wordt dat de telkunde van \cref{ch:b1:counting}, waarin elke formule
stiekem een bijectie is; voor oneindige \omterm{def:b1:logic:sets}{verzamelingen} wordt het de
weekendopgave hieronder, waar $\N$, $\Q$ en $\R$ werkelijk verschillende
groottes blijken te hebben. Datzelfde woordenboek duikt in dit volume nog
tweemaal in verfijnde vorm op: rijen (\cref{ch:b1:seq}) zijn niets anders
dan \omterm{def:b1:logic:map}{afbeeldingen} $\N \to \R$, zodat \omterm{def:b1:logic:statement}{uitspraken} over rijen \omterm{def:b1:logic:statement}{uitspraken} over
een \omterm{def:b1:logic:sets}{verzameling} \omterm{def:b1:logic:map}{afbeeldingen} zijn; en de lineaire algebra meet
vectorruimten niet met bijecties maar met \emph{lineaire} bijecties,
waarvan het bestaan door één enkel getal wordt geregeld: de dimensie
(\cref{ch:b1:findim}). Telkens als er een nieuwe ``gelijkheid'' opduikt
--- gelijkmachtigheid, isomorfie van groepen (\cref{ch:b1:structures}),
lineaire isomorfie --- herhaalt zich het patroon van
\cref{thm:b1:logic:inverse}: gelijkheid is een inverteerbare \omterm{def:b1:logic:map}{afbeelding}
die de structuur respecteert.
\end{remark}

\begin{remark}[Waar dit hoofdstuk gebruikt wordt]\label{rem:b1:logic:whereused}
Overal --- maar een paar plaatsen verdienen vermelding. De
drie-kwantoren-gymnastiek van \cref{ex:b1:logic:limit} is het dagelijks
brood van \cref{ch:b1:seq,ch:b1:continuity}: elk limietbewijs is een spel
tegen een willekeurige $\varepsilon$. \omterm{def:b1:logic:equiv}{Equivalentieklassen} keren terug als
de congruentieklassen van $\Z/n\Z$ in \cref{ch:b1:structures}, waar de
\omterm{thm:b1:logic:partition}{partitie} van \cref{thm:b1:logic:partition} een eigen algebraïsche
structuur krijgt. \omterm{def:b1:logic:order}{Ordeningsrelaties}, bovengrenzen en kleinste
bovengrenzen worden het axiomatische hart van $\R$ in
\cref{ch:b1:reals}. Injecties, surjecties en bijecties komen terug als de
lineaire \omterm{def:b1:logic:map}{afbeeldingen} van \cref{ch:b1:linmaps}, waar injectiviteit aan
één enkele vector te toetsen valt (de kern); en de weekendopgave
hieronder maakt van het kale begrip bijectie een theorie van de
\emph{groottes van oneindige \omterm{def:b1:logic:sets}{verzamelingen}}, waarvan de conclusies
(aftelbaarheid van $\Q$, overaftelbaarheid van $\R$) weer opduiken in
\cref{ch:b1:reals,ch:b1:topology}.
\end{remark}

\section{Oefeningen}

\begin{exercise}[$\star$]\label{exo:b1:logic:1}
Schrijf de negatie van elke \omterm{def:b1:logic:statement}{uitspraak}, zonder het woord ``niet'' te
gebruiken:
\begin{enumerate}
  \item $\forall x \in \R,\ \exists y \in \R,\ x + y > 0$;
  \item $\exists x \in \R,\ \forall y \in \R,\ xy = 0$;
  \item $\forall \varepsilon > 0,\ \exists \delta > 0,\ \forall x \in
        \R,\ \abs{x} \leq \delta \implies \abs{f(x)} \leq \varepsilon$
        (voor een vaste \omterm{def:b1:logic:map}{afbeelding} $f \colon \R \to \R$).
\end{enumerate}
Beslis vervolgens of de \omterm{def:b1:logic:statement}{uitspraken} (1) en (2) waar zijn.
\end{exercise}

\begin{exercise}[$\star$]\label{exo:b1:logic:2}
Zij $P, Q$ \omterm{def:b1:logic:statement}{uitspraken}. Bewijs met waarheidstabellen dat $\lnot(P
\implies Q) \iff P \land (\lnot Q)$, en leid daaruit de negatie af van:
``als een functie afleidbaar is, dan is ze continu''.
\end{exercise}

\begin{exercise}[$\star$]\label{exo:b1:logic:3}
Bewijs via contrapositie: voor $x \in \R$ geldt, als $x^3 + x \geq 2$,
dan $x \geq 1$. Bewijs vervolgens uit het ongerijmde dat er geen kleinste
strikt positief reëel getal bestaat.
\end{exercise}

\begin{exercise}[$\star$]\label{exo:b1:logic:4}
Bewijs met inductie dat voor alle $n \in \N$:
\begin{enumerate}
  \item $\sum_{k=0}^{n} 2^k = 2^{n+1} - 1$;
  \item $4^n + 5$ deelbaar is door $3$.
\end{enumerate}
\end{exercise}

\begin{exercise}[$\star$]\label{exo:b1:logic:5}
Zoek de fout in het volgende ``bewijs'' dat alle potloden dezelfde kleur
hebben. \emph{Zij $P(n)$: ``in elke \omterm{def:b1:logic:sets}{verzameling} van $n$ potloden hebben
alle potloden dezelfde kleur''. $P(1)$ is duidelijk. Neem $P(n)$ aan en
beschouw $n+1$ potloden; laat je het laatste weg, dan delen de eerste $n$
hun kleur; laat je het eerste weg, dan delen de laatste $n$ hun kleur;
dus delen alle $n+1$ hun kleur.}
\end{exercise}

\begin{exercise}[$\star$]\label{exo:b1:logic:6}
Zij $A, B, C$ deelverzamelingen van $E$. Bewijs:
\begin{enumerate}
  \item $A \setminus B = A \cap \overline{B}$;
  \item $(A \cup B) \setminus C = (A \setminus C) \cup (B \setminus
        C)$;
  \item $A \subseteq B \iff A \cup B = B \iff A \cap B = A$.
\end{enumerate}
\end{exercise}

\begin{exercise}[$\star\star$]\label{exo:b1:logic:7}
Ga voor elke \omterm{def:b1:logic:map}{afbeelding} met bewijs na of ze \omterm{def:b1:logic:inj}{injectief}, \omterm{def:b1:logic:inj}{surjectief} of
\omterm{def:b1:logic:inj}{bijectief} is:
\begin{enumerate}
  \item $f \colon \N \to \N$, $n \mapsto n + 1$;
  \item $g \colon \Z \to \Z$, $n \mapsto n + 1$;
  \item $h \colon \R \setminus \{1\} \to \R$, $x \mapsto
        \frac{x+1}{x-1}$.
\end{enumerate}
Pas voor $h$ het codomein aan zodat ze \omterm{def:b1:logic:inj}{bijectief} wordt en bereken de
inverse.
\end{exercise}

\begin{exercise}[$\star\star$]\label{exo:b1:logic:8}
Zij $f \colon E \to F$, en zij $A, A' \subseteq E$ en $B, B' \subseteq
F$.
\begin{enumerate}
  \item Bewijs $f^{-1}(B \cap B') = f^{-1}(B) \cap f^{-1}(B')$ en
        $f(A \cup A') = f(A) \cup f(A')$.
  \item Bewijs $f(A \cap A') \subseteq f(A) \cap f(A')$ en geef een
        voorbeeld waarin de inclusie strikt is.
  \item Bewijs: $f$ is \omterm{def:b1:logic:inj}{injectief} dan en slechts dan als $f(A \cap A') =
        f(A) \cap f(A')$ voor alle $A, A'$.
\end{enumerate}
\end{exercise}

\begin{exercise}[$\star\star$]\label{exo:b1:logic:9}
Zij $f \colon E \to F$ en $g \colon F \to E$ met $g \circ f =
\mathrm{id}_E$. Bewijs dat $f$ \omterm{def:b1:logic:inj}{injectief} is en $g$ \omterm{def:b1:logic:inj}{surjectief}. Geef een
voorbeeld waarin $f$ noch $g$ \omterm{def:b1:logic:inj}{bijectief} is.
\end{exercise}

\begin{exercise}[$\star\star$]\label{exo:b1:logic:10}
Definieer op $\R$: $x \mathbin{\mathcal{R}} y \iff x^2 - y^2 = x - y$.
Bewijs dat $\mathcal{R}$ een \omterm{def:b1:logic:equiv}{equivalentierelatie} is en beschrijf de
\omterm{def:b1:logic:equiv}{equivalentieklasse} van elk reëel getal $x$. Welke klassen hebben precies
één element?
\end{exercise}

\begin{exercise}[$\star\star\star$]\label{exo:b1:logic:11}
(Cantor) Zij $E$ een \omterm{def:b1:logic:sets}{verzameling}. Bewijs dat er geen surjectie van $E$ op
$\mathcal{P}(E)$ bestaat. \emph{Aanwijzing: beschouw bij een gegeven $f
\colon E \to \mathcal{P}(E)$ de \omterm{def:b1:logic:sets}{verzameling} $D = \{x \in E : x \notin
f(x)\}$.}
\end{exercise}

\begin{exercise}[$\star\star\star$]\label{exo:b1:logic:12}
Zij $f \colon E \to F$ een \omterm{def:b1:logic:map}{afbeelding}. Definieer $\Phi \colon
\mathcal{P}(F) \to \mathcal{P}(E)$ door $\Phi(B) = f^{-1}(B)$.
\begin{enumerate}
  \item Bewijs dat $f$ \omterm{def:b1:logic:inj}{surjectief} is dan en slechts dan als $\Phi$
        \omterm{def:b1:logic:inj}{injectief} is.
  \item Bewijs dat $f$ \omterm{def:b1:logic:inj}{injectief} is dan en slechts dan als $\Phi$
        \omterm{def:b1:logic:inj}{surjectief} is.
\end{enumerate}
\end{exercise}

\section{Opgave: oneindigheden vergelijken}

\begin{problem}\label{pb:b1:logic:1}
Wanneer hebben twee \omterm{def:b1:logic:sets}{verzamelingen} ``evenveel elementen''? Cantors
antwoord --- wanneer er een bijectie tussen beide bestaat --- blijkt ook
voor oneindige \omterm{def:b1:logic:sets}{verzamelingen} bruikbaar, en het splitst de oneindigheid in
werkelijk verschillende groottes. Deze opgave bouwt het volledige
gereedschap op uit de kale definities van dit hoofdstuk: de stelling van
Cantor--Schröder--Bernstein (twee injecties maken samen een bijectie),
de aftelbaarheid van $\Q$, de overaftelbaarheid van $\R$ via het
diagonaalargument, en Cantors verbluffende conclusie uit 1874:
\emph{\omterm{pb:b1:logic:1}{transcendente} getallen bestaan, en wel in overweldigende
hoeveelheid} --- zonder er ook maar één aan te wijzen.
\index{aftelbare verzameling}\index{stelling van Cantor--Schroder--Bernstein}
Schrijf overal, voor \omterm{def:b1:logic:sets}{verzamelingen} $E$ en $F$, $E \preceq F$ wanneer er
een injectie van $E$ in $F$ bestaat, en $E \approx F$ (``$E$ en $F$ zijn
\emph{gelijkmachtig}'') wanneer er een bijectie van $E$ op $F$ bestaat.

\medskip
\noindent\textbf{Deel I --- Het vocabulaire van de vergelijking.}
\begin{enumerate}
  \item Toon aan dat $\approx$ zich als een \omterm{def:b1:logic:equiv}{equivalentierelatie} gedraagt:
        $E \approx E$; als $E \approx F$, dan $F \approx E$; als $E
        \approx F$ en $F \approx G$, dan $E \approx G$. (Citeer
        nauwkeurig \cref{thm:b1:logic:inverse} en
        \cref{prop:b1:logic:comp}.)
  \item Toon aan dat $\preceq$ transitief is, en dat een injectie $f
        \colon E \to F$ altijd $E \approx f(E)$ oplevert.
  \item Zij $E \neq \emptyset$. Toon aan dat $E \preceq F$ dan en slechts
        dan als er een surjectie van $F$ op $E$ bestaat.
  \item Ga na dat $n \mapsto n + 1$ een bijectie is van $\N$ op $\N^* =
        \N \setminus \{0\}$, en dat
        \[
        \sigma(n) = \frac n2 \ \ (n \text{ even}), \qquad
        \sigma(n) = -\frac{n+1}2 \ \ (n \text{ oneven})
        \]
        een bijectie is van $\N$ op $\Z$. Een punt weglaten of naar de
        negatieve getallen verdubbelen verandert de grootte van $\N$ dus
        niet.
\end{enumerate}

\medskip
\noindent\textbf{Deel II --- De stelling van
Cantor--Schröder--Bernstein.} Zij $f \colon E \to F$ en $g \colon F \to
E$ twee injecties. Definieer
\[
C_0 = E \setminus g(F), \qquad C_{n+1} = g\bigl(f(C_n)\bigr)
\ \ (n \in \N), \qquad C = \bigcup_{n \in \N} C_n,
\]
en laat $h \colon E \to F$ elke $x \in C$ naar $f(x)$ sturen, en elke $x
\notin C$ naar de unieke $y \in F$ met $g(y) = x$.
\begin{enumerate}[resume]
  \item Ga na dat $h$ goed gedefinieerd is: als $x \notin C$, dan is $x
        \in g(F)$, en het element $y$ met $g(y) = x$ is uniek.
  \item Toon aan dat $g\bigl(f(C)\bigr) = \bigcup_{n \geq 1} C_n
        \subseteq C$. (Beelden verwisselen met verenigingen:
        \cref{exo:b1:logic:8}.)
  \item Toon aan dat $h$ \omterm{def:b1:logic:inj}{injectief} is. (Drie gevallen; toon in het
        gemengde geval $x \in C$, $x' \notin C$ aan dat $h(x) = h(x')$
        zou dwingen dat $x' \in g(f(C)) \subseteq C$.)
  \item Toon aan dat $h$ \omterm{def:b1:logic:inj}{surjectief} is: onderscheid voor gegeven $y \in
        F$ de gevallen $g(y) \notin C$ en $g(y) \in C_n$ voor zekere $n
        \geq 1$ (waarom is $g(y) \in C_0$ onmogelijk?), en wijs in elk
        geval een \omterm{def:b1:logic:map}{origineel} van $y$ aan.
  \item Besluit met de \emph{stelling van
        Cantor--Schröder--Bernstein}: als $E \preceq F$ en $F \preceq
        E$, dan $E \approx F$. Zeg in één zin wat deze \omterm{def:b1:logic:statement}{uitspraak}
        niet-triviaal maakt.
  \item Twee toepassingen. (a) Toon aan dat $\intcc01 \approx \intoo01$.
        (b) Toon aan dat $\varphi(p, q) = 2^p(2q + 1) - 1$ een bijectie
        van $\N \times \N$ op $\N$ definieert --- injectiviteit met een
        pariteitsargument, surjectiviteit met sterke inductie
        (\cref{thm:b1:logic:induction}). Bijgevolg is $\N \times \N
        \approx \N$: het rooster van gehele punten in het vlak is niet
        groter dan de lijn.
\end{enumerate}

\medskip
\noindent\textbf{Deel III --- \omterm{pb:b1:logic:1}{Aftelbare verzamelingen}.} Noem een
\omterm{def:b1:logic:sets}{verzameling} $E$ \emph{hoogstens aftelbaar} wanneer $E \preceq \N$, en
\emph{aftelbaar} wanneer $E \approx \N$.
\begin{enumerate}[resume]
  \item Toon aan dat elke oneindige deelverzameling $A \subseteq \N$
        aftelbaar is. (Definieer $\varphi(n)$ recursief als het kleinste
        element van $A \setminus \{\varphi(0), \dots, \varphi(n-1)\}$;
        toon aan dat $\varphi$ strikt stijgend is, dat $\varphi(n) \geq
        n$, en dat $\varphi$ elk element van $A$ bereikt.)
  \item Leid af dat een \omterm{def:b1:logic:sets}{verzameling} hoogstens aftelbaar is dan en slechts
        dan als ze eindig of aftelbaar is, en merk op dat vraag 9 de
        kortere weg biedt: als $E \preceq \N$ en $\N \preceq E$, dan is
        $E$ aftelbaar.
  \item Toon aan dat met $E$ en $F$ ook $E \times F$ hoogstens aftelbaar
        is. Leid af dat $\Z \times \N^*$ aftelbaar is.
  \item Toon aan dat $\Q$ aftelbaar is. (Injecteer $\Q$ in $\Z \times
        \N^*$ door elk rationaal getal onvereenvoudigbaar te schrijven
        met positieve noemer --- de uniciteit van die schrijfwijze wordt
        bewezen in \cref{ch:b1:arith}; pas dan vraag 12 toe.)
  \item Toon aan dat een aftelbare vereniging van hoogstens \omterm{pb:b1:logic:1}{aftelbare
        verzamelingen} hoogstens aftelbaar is: is elke $E_n$ ($n \in \N$)
        hoogstens aftelbaar, dan ook $\bigcup_{n \in \N} E_n$. (Stuur $x$
        naar het paar $(n, f_n(x))$ waarbij $n$ de \emph{kleinste} index
        is met $x \in E_n$.)
  \item Toon aan dat de \omterm{def:b1:logic:sets}{verzameling} van \emph{eindige}
        deelverzamelingen van $\N$ aftelbaar is. (Beeld een eindige
        deelverzameling $F$ af op $\sum_{i \in F} 2^i$; bewijs de
        injectiviteit door de grootste plaats te vergelijken waar twee
        eindige \omterm{def:b1:logic:sets}{verzamelingen} verschillen, met $\sum_{k=0}^{m-1} 2^k =
        2^m - 1$ uit \cref{exo:b1:logic:4}.)
\end{enumerate}

\medskip
\noindent\textbf{Deel IV --- Diagonalisatie.} Noteer met
$\{0,1\}^{\N}$ de \omterm{def:b1:logic:sets}{verzameling} van alle \omterm{def:b1:logic:map}{afbeeldingen} $u \colon \N \to \{0,
1\}$, dat wil zeggen de \omterm{def:b1:logic:sets}{verzameling} van binaire rijen.
\begin{enumerate}[resume]
  \item Construeer een bijectie tussen $\mathcal{P}(\N)$ en
        $\{0,1\}^{\N}$ (indicatorfuncties).
  \item (Het diagonaalargument) Zij $\Phi \colon \N \to \{0,1\}^{\N}$ een
        willekeurige \omterm{def:b1:logic:map}{afbeelding}. Beschouw de rij $d$ gegeven door $d(n) =
        1 - \Phi(n)(n)$. Toon aan dat $d$ niet in het beeld van $\Phi$
        ligt, en besluit dat $\{0,1\}^{\N}$ \emph{niet} hoogstens
        aftelbaar is. Leg in één zin uit waarom dit, via vraag 17,
        precies de stelling van Cantor (\cref{exo:b1:logic:11}) is voor
        $E = \N$.
  \item Neem aan --- zoals van de middelbare school vertrouwd, en
        streng vastgelegd in \cref{ch:b1:reals} --- dat elke $x \in
        \intco01$ precies één \emph{eigenlijke} decimale ontwikkeling $x
        = 0.d_1 d_2 d_3\dots$ heeft (een ontwikkeling die niet op een
        oneindige staart van negens eindigt). Construeer bij een
        willekeurige rij $(x_n)_{n \geq 1}$ van elementen van $\intco01$
        een $x \in \intco01$ met $x \neq x_n$ voor alle $n$: kies als
        $n$-de cijfer een $5$ wanneer het $n$-de cijfer van $x_n$ van $5$
        verschilt, en anders een $6$. Verantwoord zorgvuldig dat $x$
        eigenlijk is en elke $x_n$ ontwijkt, en besluit dat $\intco01$
        niet hoogstens aftelbaar is.
  \item Leid af dat $\R$ overaftelbaar is, en dat de \omterm{def:b1:logic:sets}{verzameling} $\R
        \setminus \Q$ van de irrationale getallen dat evenzeer is. In
        welke precieze zin zijn ``de meeste'' reële getallen irrationaal?
\end{enumerate}

\medskip
\noindent\textbf{Deel V --- Cantors stelling van 1874: \omterm{pb:b1:logic:1}{transcendente}
getallen bestaan.} Een reëel getal $x$ heet \emph{algebraïsch}
\index{algebraïsch getal} wanneer $P(x) = 0$ voor een zekere veelterm $P$
met gehele coëfficiënten die niet nul is, en \emph{transcendent}
\index{transcendent getal} in het andere geval. Neem voor dit deel aan
--- het wordt bewezen in \cref{ch:b1:poly} --- dat een veelterm van
graad $n$ die niet nul is hoogstens $n$ reële nulpunten heeft.
\begin{enumerate}[resume]
  \item Toon aan dat elk rationaal getal algebraïsch is, en geef expliciet
        veeltermen met gehele coëfficiënten die $\sqrt 2$ en $\sqrt 2 +
        \sqrt 3$ annuleren.
  \item Toon voor vaste $n \in \N$ aan dat de \omterm{def:b1:logic:sets}{verzameling} van veeltermen
        van graad hoogstens $n$ met gehele coëfficiënten aftelbaar is.
        (Injecteer haar in $\Z^{n+1}$ en pas inductie naar $n$ toe met
        vraag 13.)
  \item Leid af dat de \omterm{def:b1:logic:sets}{verzameling} van \emph{alle} veeltermen met gehele
        coëfficiënten aftelbaar is.
  \item Bewijs \emph{Cantors stelling over de algebraïsche getallen}: de
        \omterm{def:b1:logic:sets}{verzameling} $\mathcal{A}$ van de algebraïsche reële getallen is
        aftelbaar.
  \item Besluit: \omterm{pb:b1:logic:1}{transcendente} reële getallen bestaan, en de \omterm{def:b1:logic:sets}{verzameling}
        van \omterm{pb:b1:logic:1}{transcendente} getallen is overaftelbaar. Maak vervolgens in
        enkele zinnen de balans op van de hele opgave: de keten $\N
        \approx \Z \approx \Q \approx \mathcal{A}$, de strikte sprong
        naar $\R \approx$ (in wezen) $\mathcal{P}(\N)$, waar elk
        gereedschap (Cantor--Schröder--Bernstein, aftelbare
        verenigingen, de diagonaal) beslissend was --- en de filosofische
        klap van het bewijs dat er overaftelbaar veel \omterm{pb:b1:logic:1}{transcendente}
        getallen bestaan zonder er ook maar één te noemen. (Van een
        \emph{specifiek} getal als $\pi$ bewijzen dat het \omterm{pb:b1:logic:1}{transcendent}
        is, is veel moeilijker en valt buiten dit volume.)
\end{enumerate}
\end{problem}
